Let $L=\log_2 3$ and $r=\log_bx$. Writing $B=\log_2b$ and $X=\log_2x$, the first equation gives $X=5B+8$, so $B=\dfrac{8}{(r-5)}$. Substitution in the second equation and clearing denominators gives
\[
F(r)=Lr^3-12Lr^2+(37L-72)r+40-10L=0.
\]
There are three distinct real roots: since $\dfrac{3}{2}<L<2$,
\[
F(-3)=256(1-L)<0,\quad F(0)=40-10L>0,\quad
F(1)=16L-32<0,\quad F(13)=640L-936>0.
\]
Thus there is one root in each of $(-3,0),(0,1),(1,13)$.

None is excluded by a logarithm domain condition. The forbidden values are $r=0,1,5,-3,5-\dfrac{8}{L}$; the first four have nonzero $F$, and the last is not a zero because substituting $B=-L$ into the cleared logarithm equation would require $X+L=8-4L=0$. Each root therefore produces an admissible $b=2^{\dfrac{8}{(r-5)}}$ and $x=b^r$. By Vieta, their sum is $\boxed{12}$.
